\documentclass[11pt]{article}

\usepackage[margin=1in]{geometry}
\usepackage{amsmath,amssymb}
\usepackage{booktabs}
\usepackage{xcolor}
\usepackage{hyperref}
\hypersetup{
  colorlinks = true,
  linkcolor  = blue!70!black,
  citecolor  = blue!70!black,
  urlcolor   = magenta!70!black
}

\title{Optically Thin Equatorial Emission and Analytic Kerr Ray Tracing:\\
the Gralla--Lupsasca--Marrone Model}
\author{Aaron Barrios\thanks{Prepared with AI assistance; see the \hyperref[sec:ai]{AI Acknowledgments}.}}
\date{\today}

\begin{document}
\maketitle

\noindent Units: $G = c = M = 1$ throughout. Notation follows
Gralla, Lupsasca \& Marrone (GLM) \cite{GLM2020} and the Gralla--Lupsasca
geodesic papers \cite{GL2020lensing,GL2020null} on which it is built.

%----------------------------------------------------------------------
\section{Model Assumptions}

\begin{enumerate}
  \item All emission originates in the equatorial plane $\theta = \pi/2$.
  \item The emitting region is optically thin, so a light ray gains intensity
        at \emph{every} crossing of the equatorial plane. In practice this
        requires observing at $\nu \gtrsim 230$~GHz.
  \item The source is stationary and axisymmetric.
  \item Matter follows prograde circular geodesics outside the ISCO and
        plunging geodesics inside it, so emission extends down to the horizon.
  \item The finite geometric thickness of the flow is represented only by a
        multiplicative ``fudge factor'' $f$ on the lensed images.
  \item Images represent the specific intensity at a single observing
        frequency (230~GHz) for a flat emission spectrum. The overall
        normalization is fixed externally, and only the image morphology is
        predicted.
\end{enumerate}

%----------------------------------------------------------------------
\section{Spacetime}

The Kerr metric in Boyer--Lindquist coordinates is
\begin{equation}
ds^2 = -\left(1-\frac{2r}{\Sigma}\right)dt^2
       - \frac{4ar\sin^2\theta}{\Sigma}\,dt\,d\phi
       + \frac{\Sigma}{\Delta}\,dr^2 + \Sigma\,d\theta^2
       + \left(r^2+a^2+\frac{2a^2r\sin^2\theta}{\Sigma}\right)\sin^2\theta\,d\phi^2,
\end{equation}
\begin{equation}
\Sigma = r^2 + a^2\cos^2\theta, \qquad \Delta = r^2 - 2r + a^2, \qquad
r_\pm = 1 \pm \sqrt{1 - a^2},
\end{equation}
with $r_+$ the event horizon and $r_-$ the inner horizon.

%----------------------------------------------------------------------
\section{Motion of the Emitting Gas}

\subsection{\texorpdfstring{Circular orbits ($r \ge r_{\rm ISCO}$)}{Circular orbits (r >= r\_ISCO)}}

Prograde circular equatorial geodesics have
\begin{gather}
u^\mu = u^t\left(\partial_t + \Omega\,\partial_\phi\right), \qquad
\Omega = \frac{1}{r^{3/2} + a}, \qquad
u^t = \frac{r^{3/2} + a}{\sqrt{r^3 - 3r^2 + 2a\,r^{3/2}}}, \\[1ex]
E(r) = \frac{r^{3/2} - 2r^{1/2} + a}{r^{3/4}\sqrt{r^{3/2} - 3r^{1/2} + 2a}}, \qquad
L(r) = \frac{r^2 - 2a\,r^{1/2} + a^2}{r^{3/4}\sqrt{r^{3/2} - 3r^{1/2} + 2a}}.
\end{gather}
These orbits are stable for $r \ge r_{\rm ISCO}$ \cite{Bardeen1972}:
\begin{gather}
Z_1 = 1 + (1-a^2)^{1/3}\left[(1+a)^{1/3} + (1-a)^{1/3}\right], \qquad
Z_2 = \sqrt{3a^2 + Z_1^2}, \\
r_{\rm ISCO} = 3 + Z_2 - \sqrt{(3 - Z_1)(3 + Z_1 + 2Z_2)}.
\end{gather}

\subsection{\texorpdfstring{Plunging flow ($r_+ < r < r_{\rm ISCO}$)}{Plunging flow (r\_+ < r < r\_ISCO)}}

Inside the ISCO the gas falls inward along equatorial geodesics carrying the
ISCO's energy and angular momentum, with vanishing Carter constant
\cite{Cunningham1975}:
\begin{gather}
\gamma_{\rm ms} = E(r_{\rm ISCO}) = \sqrt{1 - \frac{2}{3r_{\rm ISCO}}}, \qquad
\lambda_{\rm ms} = \frac{L(r_{\rm ISCO})}{E(r_{\rm ISCO})}
= \frac{r_{\rm ISCO}^2 - 2a\,r_{\rm ISCO}^{1/2} + a^2}{r_{\rm ISCO}^{3/2} - 2r_{\rm ISCO}^{1/2} + a}, \\[1ex]
H(r) = \frac{2r - a\lambda_{\rm ms}}{\Delta(r)},
\end{gather}
\begin{align}
u^t &= \gamma_{\rm ms}\left[1 + \frac{2}{r}\bigl(1 + H\bigr)\right], \\
u^\phi &= \frac{\gamma_{\rm ms}}{r^2}\left(\lambda_{\rm ms} + aH\right), \\
u^r &= -\sqrt{\frac{2}{3r_{\rm ISCO}}}\left(\frac{r_{\rm ISCO}}{r} - 1\right)^{3/2}.
\end{align}

%----------------------------------------------------------------------
\section{Emission Profile}

The rest-frame emitted intensity $I_{\rm em}(r)$ is a free function of the
Boyer--Lindquist radius. GLM adopt a form derived from Johnson's $S_U$
distribution \cite{Johnson1949}:
\begin{equation}
I_{\rm em}(r) = J(r;\gamma,\mu,\sigma)
= \frac{\exp\!\left\{-\tfrac12\left[\gamma + \operatorname{arcsinh}\!\left(\dfrac{r-\mu}{\sigma}\right)\right]^2\right\}}{\sqrt{(r-\mu)^2 + \sigma^2}},
\label{eq:J}
\end{equation}
where $\mu$ sets the location, $\sigma$ the width, and $\gamma$ the asymmetry of
the profile. The three profiles considered in \cite{GLM2020} are listed in
Table~\ref{tab:profiles}.

\begin{table}[h]
\centering
\begin{tabular}{lcccl}
\toprule
Profile & $\gamma$ & $\mu$ & $\sigma$ & Character \\
\midrule
1 & $-3/2$ & $r_-$ & $1/2$ & peaks near the horizon, extended tail \\
2 & $0$    & $r_-$ & $1/2$ & more centrally concentrated \\
3 & $-2$   & $r_{\rm ISCO} - 1/3$ & $1/4$ & effectively truncated near the ISCO \\
\bottomrule
\end{tabular}
\caption{Emission profiles of Eq.~\eqref{eq:J} used by GLM.}
\label{tab:profiles}
\end{table}

%----------------------------------------------------------------------
\section{Photon Conserved Quantities and Screen Coordinates}

A distant observer at inclination $\theta_o$ uses Bardeen screen coordinates
$(\alpha,\beta)$ \cite{Bardeen1973}. Each pixel fixes the photon's energy-rescaled
angular momentum $\lambda$ and Carter constant $\eta$:
\begin{equation}
\lambda = -\alpha\sin\theta_o, \qquad
\eta = \beta^2 + \left(\alpha^2 - a^2\right)\cos^2\theta_o,
\end{equation}
with $\operatorname{sign}(\beta)$ the sign of $p^\theta$ at the observer.
The radial and polar potentials are
\begin{equation}
\mathcal{R}(r) = \left(r^2 + a^2 - a\lambda\right)^2 - \Delta(r)\left[\eta + (\lambda - a)^2\right],
\qquad
\Theta(\theta) = \eta + a^2\cos^2\theta - \lambda^2\cot^2\theta,
\end{equation}
and the geodesic equation, written in Mino time \cite{Mino2003}, requires the
radial and polar Mino times to agree \cite{GL2020lensing}:
\begin{equation}
I_r \equiv \int \frac{dr}{\pm\sqrt{\mathcal{R}(r)}}
= \int \frac{d\theta}{\pm\sqrt{\Theta(\theta)}} \equiv G_\theta.
\end{equation}

%----------------------------------------------------------------------
\section{Polar Motion: Mino Time to Equatorial Crossings}

Writing $u = \cos^2\theta$, the polar potential has roots
\begin{equation}
u_\pm = \Delta_\theta \pm \sqrt{\Delta_\theta^2 + \frac{\eta}{a^2}}, \qquad
\Delta_\theta = \frac12\left(1 - \frac{\eta + \lambda^2}{a^2}\right),
\end{equation}
with $u_- < 0 \le u_+$ for the ($\eta > 0$) rays that reach the equator.
The Mino time from the observer to the equatorial plane after $m$ polar
turning points is \cite{GL2020lensing,GLM2020}
\begin{equation}
G_\theta^{m} = \frac{1}{a\sqrt{-u_-}}\left[\,2m\,K\!\left(\frac{u_+}{u_-}\right) - \operatorname{sign}(\beta)\,F_o\right],
\qquad
F_o = F\!\left(\arcsin\frac{\cos\theta_o}{\sqrt{u_+}}\;\middle|\;\frac{u_+}{u_-}\right),
\label{eq:Gtheta}
\end{equation}
where $K$ and $F$ are the complete and incomplete elliptic integrals of the
first kind. Rays with $\beta < 0$ head directly toward the equator and cross
first at $m = 0$; rays with $\beta > 0$ must first reach a polar turning point,
and their first crossing is $m = 1$ (the $m = 0$ value is negative and
unphysical). The expression is singular at $a = 0$; in practice a small
floor $a \gtrsim 10^{-3}$ or the Schwarzschild limit is used.

%----------------------------------------------------------------------
\section{Radial Motion: Source Radius}

\subsection{Radial roots}

Define
\begin{equation}
\mathcal{A} = a^2 - \eta - \lambda^2, \qquad
\mathcal{B} = 2\left[\eta + (\lambda - a)^2\right], \qquad
\mathcal{C} = -a^2\eta,
\end{equation}
let $\xi_0$ be the largest real root of the resolvent cubic
\begin{equation}
\xi^3 + \mathcal{A}\,\xi^2 + \left(\frac{\mathcal{A}^2}{4} - \mathcal{C}\right)\xi - \frac{\mathcal{B}^2}{8} = 0,
\end{equation}
and set $z = \sqrt{\xi_0/2}$. The four roots of $\mathcal{R}(r)$ are
\cite{GL2020null}
\begin{equation}
r_{1,2} = -z \mp \sqrt{-\frac{\mathcal{A}}{2} - z^2 + \frac{\mathcal{B}}{4z}}, \qquad
r_{3,4} = z \mp \sqrt{-\frac{\mathcal{A}}{2} - z^2 - \frac{\mathcal{B}}{4z}},
\end{equation}
with $r_{ij} \equiv r_i - r_j$.

\subsection{Radial trajectory}

A ray traced back from $(\alpha,\beta)$ reaches radius
\cite{GL2020lensing,GLM2020}
\begin{equation}
r_s(I_r) = \frac{r_4 r_{31} - r_3 r_{41}\,\operatorname{sn}^2\!\left(\tfrac12\sqrt{r_{31}r_{42}}\,I_r - \mathcal{F}_o \,\middle|\, k\right)}
                 {r_{31} - r_{41}\,\operatorname{sn}^2\!\left(\tfrac12\sqrt{r_{31}r_{42}}\,I_r - \mathcal{F}_o \,\middle|\, k\right)},
\label{eq:rs}
\end{equation}
\begin{equation}
\mathcal{F}_o = F\!\left(\arcsin\sqrt{\frac{r_{31}}{r_{41}}}\;\middle|\;k\right),
\qquad
k = \frac{r_{32}\,r_{41}}{r_{31}\,r_{42}},
\end{equation}
after Mino time $I_r$, where $\operatorname{sn}$ is the Jacobi elliptic sine.
Note that $\mathcal{F}_o$ (built from the radial roots) differs from $F_o$ in
Eq.~\eqref{eq:Gtheta} (built from the polar roots).

\subsection{Classes of rays}

The roots are ordered $r_i \le r_j$ for $i < j$ whenever both are real, and
$r_4$ is the largest real root outside the horizon \cite{GL2020lensing}.
Rays reaching the observer fall into two classes:
\begin{itemize}
  \item \textbf{Outside the critical curve:} all four roots are real and
        $r_4 > r_+$. The ray has a radial turning point at $r_4$ and returns
        to infinity (case (2) of \cite{GL2020null}).
  \item \textbf{Inside the critical curve:} the ray has no turning point
        outside the horizon and falls into the black hole. Either all four
        roots are real with $r_4 \le r_+$ (again case (2)), or $r_3 = \bar r_4$
        are complex conjugates (case (3)).
\end{itemize}
On the critical curve itself, $r_3 = r_4$.

Equation~\eqref{eq:rs} is manifestly real in case (2), and it remains valid in
all cases \cite{GL2020null}. In case (3) it is real only after cancellation of
imaginary parts, so the following equivalent, manifestly real form is preferred
numerically. Defining
\begin{equation}
A = \sqrt{r_{32}\,r_{42}} > 0, \qquad
B = \sqrt{r_{31}\,r_{41}} > 0, \qquad
k_3 = \frac{(A+B)^2 - r_{21}^2}{4AB} \in (0,1),
\end{equation}
the source radius in case (3) is
\begin{equation}
r_s(I_r) = \frac{\left(B r_2 - A r_1\right) + \left(B r_2 + A r_1\right)\operatorname{cn}\!\left(X_3 \,\middle|\, k_3\right)}
                 {\left(B - A\right) + \left(B + A\right)\operatorname{cn}\!\left(X_3 \,\middle|\, k_3\right)},
\qquad
X_3 = \sqrt{AB}\,I_r - F\!\left(\arccos\frac{A - B}{A + B}\;\middle|\;k_3\right),
\label{eq:rs3}
\end{equation}
where $\operatorname{cn}$ is the Jacobi elliptic cosine. This follows from
inverting the case (3) radial antiderivative \cite{GL2020lensing,GL2020null}
\begin{equation}
\mathcal{I}_r^{(3)}(r) = \frac{1}{\sqrt{AB}}\,
F\!\left(\arccos\frac{A(r - r_1) - B(r - r_2)}{A(r - r_1) + B(r - r_2)}\;\middle|\;k_3\right)
\end{equation}
along the ray from the observer, $\mathcal{I}_r^{(3)}(r_s) = \mathcal{I}_r^{(3)}(\infty) - I_r$.

\subsection{Total Mino time}

The Mino time available along the portion of the ray in the Kerr exterior is
\begin{equation}
I_r^{\rm total} =
\begin{cases}
2\displaystyle\int_{r_4}^{\infty}\frac{dr}{\sqrt{\mathcal{R}(r)}}, & r_4 \in \mathbb{R},\; r_4 > r_+ \quad \text{(outside the critical curve)}, \\[2.5ex]
\displaystyle\int_{r_+}^{\infty}\frac{dr}{\sqrt{\mathcal{R}(r)}}, & \text{otherwise (inside the critical curve)}.
\end{cases}
\end{equation}
In closed form \cite{GL2020lensing}, using the case (2) antiderivative
\begin{equation}
\mathcal{I}_r^{(2)}(r) = \frac{2}{\sqrt{r_{31}r_{42}}}\,
F\!\left(\arcsin\sqrt{\frac{r - r_4}{r - r_3}\,\frac{r_{31}}{r_{41}}}\;\middle|\;k\right),
\end{equation}
\begin{equation}
I_r^{\rm total} =
\begin{cases}
\dfrac{4}{\sqrt{r_{31}r_{42}}}\,\mathcal{F}_o, & \text{outside the critical curve}, \\[2.5ex]
\dfrac{2}{\sqrt{r_{31}r_{42}}}\,\mathcal{F}_o - \mathcal{I}_r^{(2)}(r_+), & \text{inside, all roots real}, \\[2.5ex]
\dfrac{1}{\sqrt{AB}}\,F\!\left(\arccos\dfrac{A - B}{A + B}\;\middle|\;k_3\right) - \mathcal{I}_r^{(3)}(r_+), & \text{inside, } r_3 = \bar r_4.
\end{cases}
\end{equation}
For rays outside the critical curve the radial turning point at $r_4$ is
reached at $I_r = I_r^{\rm total}/2$, by symmetry of the integral.

%----------------------------------------------------------------------
\section{Equatorial Crossings}

Equating radial and polar Mino times, the $m$-th equatorial crossing of the
ray from $(\alpha,\beta)$ occurs at
\begin{equation}
r_m(\alpha,\beta) = r_s\!\left(G_\theta^{m}\right),
\qquad \text{for each } m \text{ with } 0 < G_\theta^{m} < I_r^{\rm total}.
\end{equation}
The largest admissible index is\footnote{The printed Eq.~(C8) of
\cite{GLM2020} omits the factor $I_r^{\rm total}$ in the numerator; the form
given here follows directly from the condition $G_\theta^m < I_r^{\rm total}$
with Eq.~\eqref{eq:Gtheta}.}
\begin{equation}
N(\alpha,\beta) = \left\lfloor
\frac{a\sqrt{-u_-}\,I_r^{\rm total} + \operatorname{sign}(\beta)\,F_o}{2K(u_+/u_-)}
\right\rfloor .
\end{equation}
It is convenient to label crossings by image order
\begin{equation}
n = m - H(\beta), \qquad n = 0,1,2,\dots,
\end{equation}
where $H$ is the Heaviside step function, so that $n = 0$ is the direct image
and $n \ge 1$ are the successive lensed (photon-ring) images for every pixel.

%----------------------------------------------------------------------
\section{Redshift Factor}

The redshift factor $g = \nu_{\rm obs}/\nu_{\rm em} = 1/(-p_\mu u^\mu)$, with
the photon normalized to $p_t = -1$, $p_\phi = \lambda$, depends only on the
source radius $r_s$ and on $\alpha$ (through $\lambda$) \cite{GLM2020}:
\begin{equation}
g(r_s,\alpha) =
\begin{cases}
g_{\rm orbit}, & r_s \ge r_{\rm ISCO}, \\
g_{\rm infall}, & r_s < r_{\rm ISCO}.
\end{cases}
\end{equation}

\paragraph{Circular orbits.}
\begin{equation}
g_{\rm orbit} = \frac{1}{u^t\left(1 - \Omega\lambda\right)}
= \frac{\sqrt{r_s^3 - 3r_s^2 + 2a\,r_s^{3/2}}}{r_s^{3/2} + a - \lambda}.
\end{equation}

\paragraph{Plunging flow.} With the four-velocity of Section~3.2 and the
photon's covariant radial momentum
\begin{equation}
p_r = g_{rr}\,p^r = \frac{\Sigma}{\Delta}\cdot\frac{\pm_r\sqrt{\mathcal{R}}}{\Sigma}
= \frac{\pm_r\sqrt{\mathcal{R}(r_s)}}{\Delta(r_s)},
\end{equation}
the redshift factor is\footnote{The printed Eq.~(B15) of \cite{GLM2020}
multiplies by $\Delta(r_s)$ rather than dividing. That form is dimensionally
inconsistent ($u^t$ is dimensionless while $\Delta\sqrt{\mathcal{R}}$ scales
as $r^4$); the expression here follows from $p_r = g_{rr}p^r$.}
\begin{equation}
g_{\rm infall} = \frac{1}{\,u^t - \lambda\,u^\phi - u^r\,\dfrac{\pm_r\sqrt{\mathcal{R}(r_s)}}{\Delta(r_s)}\,}.
\end{equation}
Here $\pm_r$ is the sign of the photon's radial motion at emission: $+$ for
rays inside the critical curve, and for rays outside it $+$ if the crossing
lies on the outgoing leg after the turning point ($I_r < I_r^{\rm total}/2$),
$-$ otherwise.

This single factor contains the Doppler shift, the transverse Doppler effect,
and gravitational redshift; none of them should be applied separately.

%----------------------------------------------------------------------
\section{Observed Image}

Tracing each ray back through the emitting plane and summing over its
crossings gives \cite{GLM2020}
\begin{equation}
I_{\rm obs}(\alpha,\beta) = \sum_{n} f_n\; g^p\!\left(r_n,\alpha\right)\, I_{\rm em}(r_n),
\qquad
f_n = \begin{cases} 1, & n = 0, \\ f, & n \ge 1, \end{cases}
\qquad p = 3,
\label{eq:Iobs}
\end{equation}
where the sum runs over all admissible crossings of Section~8.

\paragraph{Power of the redshift factor.} Light emitted at frequency
$\nu/g$ is observed at $\nu$, and invariance of $I_\nu/\nu^3$ along the ray
gives
\begin{equation}
I_\nu^{\rm obs}(\nu) = g^3\, I_\nu^{\rm em}(\nu/g).
\end{equation}
For an emission spectrum $I_\nu^{\rm em} \propto \nu^{-\alpha_\nu}$ this
becomes $I_\nu^{\rm obs} = g^{3+\alpha_\nu} I_\nu^{\rm em}(\nu)$, so
$p = 3 + \alpha_\nu$; integrating over frequency instead
($d\nu_{\rm obs} = g\,d\nu_{\rm em}$) gives the bolometric value $p = 4$, which
is the choice made by GLM \cite{GLM2020}. Our images represent EHT
observations at a single frequency, 230~GHz, where the spectra of Sgr~A$^*$
and M87$^*$ are close to flat ($\alpha_\nu \approx 0$), so we adopt $p = 3$.
A larger $p$ strengthens the Doppler asymmetry, since the brightness contrast
scales as $(g_{\max}/g_{\min})^p$, and further dims the strongly redshifted
inner region; GLM find the difference between $p = 3$ and $p = 4$ small and
largely degenerate with the choice of emission profile.

\paragraph{Fudge factor.} The fudge factor $f$ mimics the extra path length that lensed rays accumulate
through a geometrically thick flow: $f = 1$ is a strictly equatorial disk,
and GLM survey $f \in \{1, 1.5, 2\}$ with $f = 1.5$ as their best guess.

The fudge factor is assigned by image order $n$, so it multiplies the lensed
images $n \ge 1$ across the entire screen and never the direct image, as in
the AART implementation \cite{AART2023}. GLM write the same rule as $f_m$,
where in their Eq.~(6) $m$ counts equatorial crossings along the backward ray
(our $n$); in their Appendix~C the same symbol denotes the number of polar
turning points, as in Eq.~\eqref{eq:Gtheta}. The two coincide only for
$\beta < 0$, which is why this document uses $n$ for crossing order and $m$
for turning points.

Since the normalization of $I_{\rm em}$ is arbitrary, the normalized image
\begin{equation}
\hat I(\alpha,\beta) = \frac{I_{\rm obs}(\alpha,\beta)}{\max I_{\rm obs}}
\end{equation}
depends on
\begin{equation}
\underbrace{(a,\;\theta_o)}_{\text{physical parameters}}
\quad\text{and}\quad
\underbrace{(\gamma,\;\mu,\;\sigma,\;f)}_{\text{emission nuisance parameters}}.
\end{equation}
The physical scale of the image follows from $M/D$; for M87$^*$,
$M/D \approx 3.62\,\psi~\mu{\rm as}$ with $\psi = 1$ corresponding to
$M = 6.2\times10^9\,M_\odot$ and $D = 16.9$~Mpc \cite{GLM2020}.

%----------------------------------------------------------------------
\section{Parameter Choices in GLM}

\begin{table}[h]
\centering
\begin{tabular}{lccc}
\toprule
 & Survey & Best guess & Alternative \\
\midrule
Spin $a$             & $0.01,\;0.5,\;0.99$       & $0.94$ & $0.01$ \\
Inclination $\theta_o$ & $10^\circ,\;20^\circ,\;30^\circ$ & $17^\circ$ & $10^\circ$ \\
Emission profile     & 1, 2, 3                   & 1      & truncated at the ISCO \\
Fudge factor $f$     & $1,\;1.5,\;2$             & $1.5$  & $1$ \\
\bottomrule
\end{tabular}
\caption{Parameters used by GLM for M87$^*$. Inclinations above $30^\circ$
are disfavored by observations of M87$^*$.}
\end{table}

%----------------------------------------------------------------------
\section{Limitations and Validation}

\subsection{Limitations}

\begin{itemize}
  \item \textbf{Equatorial approximation.} Real low-luminosity flows are
        geometrically thick; $f$ is a phenomenological correction, not a
        derived quantity.
  \item \textbf{Optical depth.} Synchrotron self-absorption is neglected; the
        model is valid only where the flow is optically thin
        ($\nu \gtrsim 230$~GHz for M87$^*$ and Sgr~A$^*$).
  \item \textbf{Stationarity and axisymmetry.} Turbulent variability is
        absent from the base model.
  \item \textbf{Spectral shape.} A flat spectrum ($\alpha_\nu = 0$, hence
        $p = 3$) is assumed everywhere; in a real flow the spectral index
        varies with radius, which would change the effective $p$ across the
        image.
  \item \textbf{Spin information.} Because emission continues inside the ISCO,
        there is no sharp inner edge at $r_{\rm ISCO}$; spin enters through
        the plunging-region redshift, the Doppler asymmetry, and the shape and
        size of the $n \ge 1$ photon rings.
\end{itemize}

\subsection{Validation tests}

Several conventions above are not stated explicitly in \cite{GLM2020} and
must be confirmed numerically: the image-order labels $n = m - H(\beta)$ and
the assignment of $f$ by image order, the sign $\pm_r$ in $g_{\rm infall}$,
and the agreement of Eqs.~\eqref{eq:rs} and \eqref{eq:rs3} in case (3). A
generator implementing this model should pass the following tests before it
is used to build a training set.

\begin{enumerate}
  \item \textbf{Schwarzschild limit ($a \to 0$).} The image must reduce to the
        standard non-spinning result, with the $n \ge 1$ rings converging to
        the circle $\sqrt{\alpha^2 + \beta^2} = 3\sqrt{3} \approx 5.196$.
        \emph{Tests:} the radial roots, Eq.~\eqref{eq:rs}, and the crossing
        count $N$.
  \item \textbf{Critical curve.} For $a > 0$, the $n \ge 1$ rings must
        converge to Bardeen's analytic critical curve \cite{Bardeen1973},
        traced by the bound photon orbits $\tilde r \in [\tilde r_-, \tilde r_+]$:
        \begin{equation}
        \lambda(\tilde r) = a + \frac{\tilde r}{a}\left[\tilde r - \frac{2\Delta(\tilde r)}{\tilde r - 1}\right],
        \qquad
        \eta(\tilde r) = \frac{\tilde r^3}{a^2}\left[\frac{4\Delta(\tilde r)}{(\tilde r - 1)^2} - \tilde r\right],
        \end{equation}
        \begin{equation}
        \tilde r_\pm = 2\left[1 + \cos\!\left(\frac23\arccos(\pm a)\right)\right],
        \end{equation}
        mapped to the screen with the relations of Section~5.
        \emph{Tests:} the image-order labels and the treatment of rays inside
        versus outside the critical curve.
  \item \textbf{Face-on symmetry ($\theta_o \to 0$).} The image must be
        axisymmetric about the origin.
        \emph{Tests:} the sign conventions for $\beta$ and in $G_\theta^m$.
  \item \textbf{Doppler asymmetry.} For $\theta_o > 0$, the side of the image
        where the gas approaches the observer must be brighter.
        \emph{Tests:} the sign of $\lambda$ and the redshift factor $g$.
  \item \textbf{Continuity at the ISCO.} The observed intensity must be
        continuous across the image of $r_{\rm ISCO}$, where $g_{\rm orbit}$
        hands over to $g_{\rm infall}$.
        \emph{Tests:} the plunging four-velocity, the factor $\Delta$, and the
        sign $\pm_r$ in $g_{\rm infall}$.
  \item \textbf{Comparison with AART.} Images and $r_m(\alpha,\beta)$ maps must
        agree pixel by pixel with the AART code \cite{AART2023}, which
        implements this model.
        \emph{Tests:} the full pipeline.
\end{enumerate}

%----------------------------------------------------------------------
\section*{AI Acknowledgments}
\phantomsection\label{sec:ai}

This document was prepared with the assistance of Claude (Anthropic), used
through Claude Code. The author directed the scope, content, and modeling
choices. Claude drafted the text and equations from the cited literature,
checked them against the original papers and the AART source code, verified
the references, and supplied derivations and validation tests where the
literature leaves details unstated. The author is responsible for the final
content.

\begin{thebibliography}{9}

\bibitem{GLM2020}
S.~E. Gralla, A.~Lupsasca, and D.~P. Marrone,
``The shape of the black hole photon ring: A precise test of strong-field
general relativity,''
\emph{Phys.\ Rev.\ D} \textbf{102}, 124004 (2020), arXiv:2008.03879.

\bibitem{GL2020lensing}
S.~E. Gralla and A.~Lupsasca,
``Lensing by Kerr black holes,''
\emph{Phys.\ Rev.\ D} \textbf{101}, 044031 (2020), arXiv:1910.12873.

\bibitem{GL2020null}
S.~E. Gralla and A.~Lupsasca,
``Null geodesics of the Kerr exterior,''
\emph{Phys.\ Rev.\ D} \textbf{101}, 044032 (2020), arXiv:1910.12881.

\bibitem{Bardeen1972}
J.~M. Bardeen, W.~H. Press, and S.~A. Teukolsky,
``Rotating black holes: Locally nonrotating frames, energy extraction, and
scalar synchrotron radiation,''
\emph{Astrophys.\ J.} \textbf{178}, 347--370 (1972).

\bibitem{Bardeen1973}
J.~M. Bardeen, ``Timelike and null geodesics in the Kerr metric,''
in \emph{Black Holes (Les Astres Occlus)},
ed.\ C.~DeWitt and B.~S. DeWitt (Gordon and Breach, New York, 1973),
pp.~215--239.

\bibitem{Cunningham1975}
C.~T. Cunningham,
``The effects of redshifts and focusing on the spectrum of an accretion disk
around a Kerr black hole,''
\emph{Astrophys.\ J.} \textbf{202}, 788--802 (1975).

\bibitem{Mino2003}
Y.~Mino,
``Perturbative approach to an orbital evolution around a supermassive black
hole,''
\emph{Phys.\ Rev.\ D} \textbf{67}, 084027 (2003).

\bibitem{Johnson1949}
N.~L. Johnson,
``Systems of frequency curves generated by methods of translation,''
\emph{Biometrika} \textbf{36}, 149--176 (1949).

\bibitem{AART2023}
A.~C\'ardenas-Avenda\~no, A.~Lupsasca, and H.~Zhu,
``Adaptive analytical ray tracing of black hole photon rings,''
\emph{Phys.\ Rev.\ D} \textbf{107}, 043030 (2023), arXiv:2211.07469.

\end{thebibliography}

\end{document}
